\section{The rectangle graph and polynomial conditioning}\label{sec:graph}
Here $m\ge2$ need not be odd. For a cyclic $\GH(m,2)$ seed, put a vertex at each cell $(a,j)$ and join $(a,j)$ to $(b,\ell)$ when
\begin{equation}\label{eq:edge}
 a\ne b,\quad j\ne\ell,\quad G_{aj}G_{b\ell}=G_{a\ell}G_{bj}.
\end{equation}
Each row-ratio class has two columns, so the other diagonal $(a,\ell),(b,j)$ is also joined. The column GH counts follow from all-power column orthogonality and Fourier inversion. Thus every vertex has exactly one neighbor in each other row and exactly one in each other column: the graph is simple and $(n-1)$-regular.

\begin{lemma}\label{lem:componentweight}
Every component $C$ contains the same number $c_C$ of cells in each row and each column. It has $nc_C$ vertices, and $\sum_Cc_C=n$.
\end{lemma}
\begin{proof}
The neighbor correspondence between any two rows is a bijection of their cells. It restricts to a bijection on cells belonging to any fixed component, since edges cannot leave it. The column argument is identical, and total-size counting identifies the row and column constants.
\end{proof}

\subsection{The component formula for the full kernel}
Write $V_G$ for the simultaneous dephased first-order kernel. Real phase directions are taken modulo row and column phase directions. Each class has a unique double-centered representative $X$, with every row and column sum zero, obtained by subtracting row and column means and adding the global mean. This preserves the tangent equations because phase gauges are exact symmetries.

\begin{theorem}\label{thm:graphkernel}
A double-centered real matrix is in the simultaneous first-order kernel exactly when it is constant on every rectangle-graph component. Consequently, if the graph has $r$ components,
\begin{equation}\label{eq:dimension}
 \dim V_G=r-1.
\end{equation}
In particular $L_G$ has full column rank exactly when the graph is connected.
\end{theorem}
\begin{proof}
Finite Fourier inversion of the first-order row equations says that the sums of $X_{aj}-X_{bj}$ over the two-element ratio classes are equal. Their total is zero by double centering, so each is zero. The corresponding column assertion holds as well: differentiating square unitarity shows that a vanishing row Gram derivative at an exact square Hadamard matrix implies a vanishing column Gram derivative.

For a rectangle in~\eqref{eq:edge}, the row and column equations are
\[
 X_{aj}+X_{a\ell}-X_{bj}-X_{b\ell}=0,\qquad
 X_{aj}+X_{bj}-X_{a\ell}-X_{b\ell}=0.
\]
Adding and subtracting gives $X_{aj}=X_{b\ell}$ and $X_{a\ell}=X_{bj}$. Thus the values are component-constant. Conversely diagonal constancy gives these equations, hence all row derivatives by Fourier inversion.

If the component values are $u_C$, Lemma~\ref{lem:componentweight} makes every row and column sum equal to $\sum_Cc_Cu_C$. Double centering imposes precisely this one independent equation, proving~\eqref{eq:dimension}. Dephased and double-centered spaces are complements to the same $(2n-1)$-dimensional gauge space.
\end{proof}
For a proper component indicator $W$, the matrix $W-(c_C/n)J$ is double-centered and its dephased representative is
\[
 X_{ij}=W_{ij}-W_{i0}-W_{0j}+W_{00}.
\]
Thus a disconnected graph supplies rational kernel directions without solving a cyclotomic linear system.

\subsection{Edge estimates and graph diameter}
Suppose the graph is connected, of diameter $D$. Write $K=G\circ e^{iX}$ in small dephased principal phases, let $Y$ be its double centering, and put $y=\norm Y_\mx$. The corresponding phase-gauged $K'=G\circ e^{iY}$ has the same residual.

\begin{lemma}\label{lem:edgebound}
For every edge $(a,j)\sim(b,\ell)$,
\begin{equation}\label{eq:edgeestimate}
 |Y_{aj}-Y_{b\ell}|\le\frac hm\eps+n(m-1)y^2.
\end{equation}
\end{lemma}
\begin{proof}
Taylor expansion of the power-$k$ row correlation gives
\[
 \abs{\sum_t(G_{at}/G_{bt})^k(Y_{at}-Y_{bt})}
 \le\eps/k+2nky^2.
\]
The zero-frequency sum vanishes by double centering. Fourier inversion for the two-element class $j,\ell$ yields
\[
 |Y_{aj}+Y_{a\ell}-Y_{bj}-Y_{b\ell}|
 \le(h/m)\eps+n(m-1)y^2.
\]
By~\eqref{eq:column-residual}, the column-pair sum has the same bound. Add the two sums and divide by two to obtain~\eqref{eq:edgeestimate}.
\end{proof}
Since $Y$ has mean zero, its maximum absolute value is at most its range. Graph paths therefore give
\[
 y\le D[(h/m)\eps+n(m-1)y^2].
\]
If $y\le1/[2Dn(m-1)]$, absorb the quadratic term to get $y\le2Dh\eps/m$. Recovering dephased phases by $X_{ij}=Y_{ij}-Y_{i0}-Y_{0j}+Y_{00}$ proves
\begin{equation}\label{eq:graphlinear}
 \norm{K-G}_\mx\le\norm X_\mx\le\frac{8Dh}{m}\eps.
\end{equation}

For completeness, a connected graph on $n^2$ vertices of minimum degree $n-1$ has diameter at most $3n-1$. On a geodesic of length $D$, select positions $0,3,6,\ldots,3\lfloor D/3\rfloor$. Their closed neighborhoods are disjoint, since overlap would connect two selected vertices within distance two. Each neighborhood has at least $n$ vertices, so
\[
 [\lfloor D/3\rfloor+1]n\le n^2,\qquad D\le3n-1<6m.
\]
Applied inside a component of row weight $c_C$, the same argument gives component diameter at most $3c_C-1$.

\subsection{An explicit residual-only entrance condition}
Now take odd $m\ge3$ and round by Theorem~\ref{thm:odd}. The principal phases satisfy
\[
 \norm X_\mx\le1024\sqrt{\eps/m},\qquad
 y\le4\norm X_\mx\le4096\sqrt{\eps/m}.
\]
For the rounded seed's graph, the local radius follows from
\begin{equation}\label{eq:graphentrance}
 \eps\le\frac1{2^{28}D^2m(m-1)^2}.
\end{equation}
Since $D<6m$, the simpler sufficient restriction is $\eps\le2^{-34}m^{-5}$, already below the global threshold.

\begin{theorem}[Polynomial linear certificate]\label{thm:polynomial}
For odd $m\ge3$, let $H$ have exactly unit-modulus entries and residual at most $2^{-34}m^{-5}$. If the rectangle graph of its rounded dephased seed $G$ is connected, then
\[
 \norm{K-G}_\mx\le48H_{m-1}\eps.
\]
\end{theorem}
The graph is computable by integer modular comparisons and traversal. The theorem is a per-seed certificate; it is not a proof that every odd seed's graph is connected. A disconnected odd seed would instead yield a genuine square-root obstruction by Proposition~\ref{prop:odd-dichotomy}. Discreteness of the exact set alone does not imply a zero tangent kernel.

For the two classical seeds of Section~\ref{sec:integer}, exact traversal gives one component and diameters $3$ and $4$. Formula~\eqref{eq:graphlinear} improves their local constants to $12$ and $40/3$, respectively. These are conditioning improvements, not new isolation results.
